% Part: computability
% Chapter: computability-theory
% Section: s-m-n

\documentclass[../../../include/open-logic-section]{subfiles}

\begin{document}

\olfileid{cmp}{thy}{smn}
\olsection{$s$-$m$-$n$ પ્રમેય}

\begin{explain}
આગળનું પ્રમેય ``$s$-$m$-$n$ પ્રમેય'' તરીકે ઓળખાય છે; તેનું કારણ
થોડી વારમાં સ્પષ્ટ થશે. અઘરું કામ પ્રમેય ખરેખર શું કહે છે તે
સમજવાનું છે; વિધાન સમજાઈ જાય પછી તે ઘણું સ્વાભાવિક લાગશે.
\end{explain}

\begin{thm}
\ollabel{thm:s-m-n}
  પ્રાકૃતિક સંખ્યાઓની દરેક જોડી $n$ અને~$m$ માટે એવું આદિમ
  પુનરાવર્તી વિધેય~$s^m_n$ છે કે દરેક અનુક્રમ
  $e$, $a_0$, \dots, $a_{m-1}$, $y_0$, \dots, $y_{n-1}$ માટે
  \[
  \cfind{s^m_n(e, a_0, \dots, a_{m-1})}[n](y_0, \dots, y_{n-1}) \simeq
  \cfind{e}[m+n](a_0, \dots, a_{m-1}, y_0, \dots, y_{n-1}).
\]
\end{thm}

\begin{explain}
$s^m_n$ને \emph{પ્રોગ્રામો} પર કાર્ય કરતું માનવાથી મદદ મળે છે.
એટલે કે, $s^m_n$ને $(m+n)$-આર્ગ્યુમેન્ટવાળા વિધેયનો
પ્રોગ્રામ~$e$ અને નિશ્ચિત ઇનપુટ $a_0$, \dots, $a_{m-1}$ મળે;
તે બાકીના આર્ગ્યુમેન્ટોના $n$-આર્ગ્યુમેન્ટવાળા વિધેય માટે
$s^m_n(e,a_0,\dots,a_{m-1})$ પ્રોગ્રામ આપે છે.
\sourcecorrection{OLCT-002}{સ્થિર અંગ્રેજી મૂળ આ છેલ્લા
પ્રોગ્રામના સૂચકમાં અને ટ્યુરિંગ-મશીનના વર્ણનમાં $e$ને બદલે $x$
લખે છે. ઉપરના પ્રમેય અને તે જ વાક્યના પ્રોગ્રામ~$e$ મુજબ બંને
જગ્યાએ $e$ રાખ્યું છે.}
\iftag{TMs}{ જો $e$ને ટ્યુરિંગ મશીનનું વર્ણન માનીએ, તો
$s^m_n(e,a_0,\dots,a_{m-1})$ એ એવું ટ્યુરિંગ મશીન છે જે
$y_0$, \dots, $y_{n-1}$ ઇનપુટ પર ઇનપુટ-શબ્દમાળાની શરૂઆતમાં
$a_0$, \dots, $a_{m-1}$ જોડે છે અને પછી~$e$ ચલાવે છે.
દરેક $s^m_n$ એ યોગ્ય ટ્યુરિંગ મશીનનો કોડ શોધતું આદિમ
પુનરાવર્તી વિધેય છે.}{}
\end{explain}

\end{document}
